% GENERATED FILE. Edit canonical lessons, then run npm run generate:books.
% canonical-source-hash: fnv1a64:18801fbd2287c2c9
% canonical-lessons: TE-C82-sixth, TE-C82-seventh, TE-C82-eighth, TE-C82-ninth, TE-C82-tenth, TE-C82-eleventh

\chapter{Sixth to Eleventh}
\label{ch:te-ordinals-to-tenth}
\hlchaptermodality{82}

\begin{chapteropening}
\textbf{By the end of this chapter:} \emph{I can count in order all the way to tenth in Telugu and beyond it, because the ending that makes an ordinal does not stop at ten.}

The ending taken past ten, which shows it was never a fact about the numbers one to ten but about numbers.
\end{chapteropening}

\section[ārava]{\te{ఆరవ} (ārava) — sixth}

\label{lesson:TE-C82-sixth}



\begin{quote}
Recalls open the chapter, at the four measured distances.

\begin{itemize}
\raggedright
  \item \emph{Recall:} say \emph{aidava} — \textbf{R1}, one lesson back, before it decays
  \item \emph{Recall:} say \emph{modaṭi} — \textbf{R2}, the first real retrieval, five to fifteen lessons back
  \item \emph{Recall:} say \emph{veḷḷavaccu} — \textbf{R2}, a second word from the same distance
  \item \emph{Recall:} say \emph{eluka} — \textbf{R3}, durable, twenty to sixty lessons back
  \item \emph{Recall:} say \emph{ḍabbu} — \textbf{R3}, a second word from the same distance
  \item \emph{Recall:} say \emph{arthaṁ cēsukō} — \textbf{R4}, recognition at distance, eighty lessons back or more
  \item \emph{Recall:} say \emph{the letter tha} — \textbf{R4}, a second word from the same distance
\end{itemize}
\end{quote}

\subsection*{You'll want to know: \te{ఆరవ}}
\textbf{\te{ఆరవ}} (\emph{ārava}) — \textquotedblleft{}\textbf{sixth}\textquotedblright{}.

\textbf{\te{ఆరు}} $\to$ \textbf{\te{ఆరవ}}, and this is the shortest ordinal in the ten: three syllables, two of them the number itself.

\textbf{\te{ఆరవ} \te{రోజు}} (\emph{ārava rōju}) — the sixth day.

The numbers six to ten were the ones the first numbers lesson called Telugu's wandering half, the place where it drifts furthest from Tamil and Kannada. The ordinal ending does not care. Whatever shape the cardinal arrived in, \textbf{-\te{వ}} goes on it.

\subsection*{Your turn}
Say these aloud:

\begin{itemize}
\raggedright
  \item \emph{ārava}
  \item the swap — \emph{āru} $\to$ \emph{ārava}
  \item \emph{ārava rōju} — the sixth day
\end{itemize}

\begin{tcolorbox}[breakable,colback=teal!4,colframe=teal!35!black,title={Before you move on}]
Say the sixth. (\textbf{\te{ఆరవ}}.) What was traded for the ending? (\textbf{The final \emph{u} of \te{ఆరు}}.) And does the ending change for the wandering half of the numbers? (\textbf{No} — it is the same \emph{-\te{వ}}.)
\end{tcolorbox}
\section[ēḍava]{\te{ఏడవ} (ēḍava) — seventh}

\label{lesson:TE-C82-seventh}



\begin{quote}
Recalls first, then the seventh.

\begin{itemize}
\raggedright
  \item \emph{Recall:} say \emph{ārava} — \textbf{R1}, one lesson back, before it decays
  \item \emph{Recall:} say \emph{reṇḍava} — \textbf{R2}, the first real retrieval, five to fifteen lessons back
  \item \emph{Recall:} say \emph{lōpaliki raṇḍi} — \textbf{R2}, a second word from the same distance
  \item \emph{Recall:} say \emph{kappa} — \textbf{R3}, durable, twenty to sixty lessons back
  \item \emph{Recall:} say \emph{dhara} — \textbf{R3}, a second word from the same distance
  \item \emph{Recall:} say \emph{sahāyaṁ cēyu} — \textbf{R4}, recognition at distance, eighty lessons back or more
  \item \emph{Recall:} say \emph{nāku telugu iṣṭaṁ} — \textbf{R4}, a second word from the same distance
\end{itemize}
\end{quote}

\subsection*{You'll want to know: \te{ఏడవ}}
\textbf{\te{ఏడవ}} (\emph{ēḍava}) — \textquotedblleft{}\textbf{seventh}\textquotedblright{}.

\textbf{\te{ఏడు}} $\to$ \textbf{\te{ఏడవ}}. Hold the \textbf{ē} long — the letter \textbf{\te{ఏ}} is the long one, and its short partner \textbf{\te{ఎ}} is a different letter and a different word.

\textbf{\te{ఏడవ} \te{వారం}} (\emph{ēḍava vāraṁ}) — the seventh week.

Seven weeks is a useful thing to be able to say and a good place to hear that the ordinal is doing work no cardinal can do: \emph{\te{ఏడు} \te{వారం}} would be seven weeks, a quantity. \emph{\te{ఏడవ} \te{వారం}} is the seventh one, a position.

\subsection*{Your turn}
Say these aloud:

\begin{itemize}
\raggedright
  \item \emph{ēḍava}
  \item the swap — \emph{ēḍu} $\to$ \emph{ēḍava}
  \item \emph{ēḍu vāraṁ} is seven weeks; \emph{ēḍava vāraṁ} is the seventh
\end{itemize}

\begin{tcolorbox}[breakable,colback=teal!4,colframe=teal!35!black,title={Before you move on}]
Say the seventh. (\textbf{\te{ఏడవ}}.) Which vowel opens it, the long or the short one? (\textbf{The long \te{ఏ}}.) And what is the difference between \emph{\te{ఏడు} \te{వారం}} and \emph{\te{ఏడవ} \te{వారం}}? (\textbf{Seven weeks} against \textbf{the seventh week}.)
\end{tcolorbox}
\section[enimidava]{\te{ఎనిమిదవ} (enimidava) — eighth}

\label{lesson:TE-C82-eighth}



\begin{quote}
Recalls first, then the eighth.

\begin{itemize}
\raggedright
  \item \emph{Recall:} say \emph{ēḍava} — \textbf{R1}, one lesson back, before it decays
  \item \emph{Recall:} say \emph{mūḍava} — \textbf{R2}, the first real retrieval, five to fifteen lessons back
  \item \emph{Recall:} say \emph{śubhākāṅkṣalu} — \textbf{R2}, a second word from the same distance
  \item \emph{Recall:} say \emph{dōma} — \textbf{R3}, durable, twenty to sixty lessons back
  \item \emph{Recall:} say \emph{tūkaṁ} — \textbf{R3}, a second word from the same distance
  \item \emph{Recall:} say \emph{dayacēsi ṭī} — \textbf{R4}, recognition at distance, eighty lessons back or more
  \item \emph{Recall:} say \emph{the vowel sign ṛ} — \textbf{R4}, a second word from the same distance
\end{itemize}
\end{quote}

\subsection*{You'll want to know: \te{ఎనిమిదవ}}
\textbf{\te{ఎనిమిదవ}} (\emph{enimidava}) — \textquotedblleft{}\textbf{eighth}\textquotedblright{}.

\textbf{\te{ఎనిమిది}} $\to$ \textbf{\te{ఎనిమిదవ}}, and here the vowel that goes is an \textbf{i}, not a \textbf{u}. Every ordinal so far has traded a final \emph{u}; this one and the next trade a final \emph{i}, because that is what their cardinals happen to end in.

Which is the useful way to hold the rule: it is not \textquotedblleft{}drop the \emph{u}\textquotedblright{}, it is \textbf{drop whatever vowel is at the end}. The ending does not know which vowel it replaced.

Note the short \textbf{\te{ఎ}} at the front, against the long \textbf{\te{ఏ}} of the last lesson.

\subsection*{Your turn}
Say these aloud:

\begin{itemize}
\raggedright
  \item \emph{enimidava}
  \item the swap — \emph{enimidi} $\to$ \emph{enimidava}, an \emph{i} this time
  \item short \emph{\te{ఎ}} here, long \emph{\te{ఏ}} in \emph{ēḍava}
\end{itemize}

\begin{tcolorbox}[breakable,colback=teal!4,colframe=teal!35!black,title={Before you move on}]
Say the eighth. (\textbf{\te{ఎనిమిదవ}}.) Which vowel came off this time? (\textbf{An \emph{i}}, not a \emph{u}.) And how is the rule best held? (\textbf{Drop whatever final vowel is there}.)
\end{tcolorbox}
\section[tommidava]{\te{తొమ్మిదవ} (tommidava) — ninth}

\label{lesson:TE-C82-ninth}



\begin{quote}
Recalls first, then the ninth.

\begin{itemize}
\raggedright
  \item \emph{Recall:} say \emph{enimidava} — \textbf{R1}, one lesson back, before it decays
  \item \emph{Recall:} say \emph{nālugava} — \textbf{R2}, the first real retrieval, five to fifteen lessons back
  \item \emph{Recall:} say \emph{uḍuta} — \textbf{R3}, durable, twenty to sixty lessons back
  \item \emph{Recall:} say \emph{bēraṁ} — \textbf{R3}, a second word from the same distance
  \item \emph{Recall:} say \emph{idi} — \textbf{R4}, recognition at distance, eighty lessons back or more
  \item \emph{Recall:} say \emph{adi} — \textbf{R4}, a second word from the same distance
\end{itemize}
\end{quote}

\subsection*{You'll want to know: \te{తొమ్మిదవ}}
\textbf{\te{తొమ్మిదవ}} (\emph{tommidava}) — \textquotedblleft{}\textbf{ninth}\textquotedblright{}.

\textbf{\te{తొమ్మిది}} $\to$ \textbf{\te{తొమ్మిదవ}}. The doubled \textbf{\te{మ్మ}} in the middle is held, exactly as it is in the cardinal — the ending is added at the far end of the word and touches nothing on its way.

\textbf{\te{తొమ్మిదవ} \te{రోజు}} (\emph{tommidava rōju}) — the ninth day.

Second \emph{i}-ending in a row, and the same trade. One more and the ten are closed.

\subsection*{Your turn}
Say these aloud:

\begin{itemize}
\raggedright
  \item \emph{tommidava} — hold the \emph{mm}
  \item the swap — \emph{tommidi} $\to$ \emph{tommidava}
  \item \emph{tommidava rōju} — the ninth day
\end{itemize}

\begin{tcolorbox}[breakable,colback=teal!4,colframe=teal!35!black,title={Before you move on}]
Say the ninth. (\textbf{\te{తొమ్మిదవ}}.) What happens to the doubled \emph{mm}? (\textbf{Nothing} — it is held, as in the cardinal.) And which vowel came off? (\textbf{The final \emph{i}}.)
\end{tcolorbox}
\section[padava]{\te{పదవ} (padava) — tenth}

\label{lesson:TE-C82-tenth}



\begin{quote}
Recalls first, then the tenth, which closes the series.

\begin{itemize}
\raggedright
  \item \emph{Recall:} say \emph{tommidava} — \textbf{R1}, one lesson back, before it decays
  \item \emph{Recall:} say \emph{aidava} — \textbf{R2}, the first real retrieval, five to fifteen lessons back
  \item \emph{Recall:} say \emph{sālīḍu} — \textbf{R3}, durable, twenty to sixty lessons back
  \item \emph{Recall:} say \emph{tīpi} — \textbf{R3}, a second word from the same distance
  \item \emph{Recall:} say \emph{akkaḍa} — \textbf{R4}, recognition at distance, eighty lessons back or more
  \item \emph{Recall:} say \emph{evaru} — \textbf{R4}, a second word from the same distance
\end{itemize}
\end{quote}

\subsection*{You'll want to know: \te{పదవ}}
\textbf{\te{పదవ}} (\emph{padava}) — \textquotedblleft{}\textbf{tenth}\textquotedblright{}.

\textbf{\te{పది}} $\to$ \textbf{\te{పదవ}}. Two syllables become three, and the trade is the same one you have now made nine times.

\textbf{\te{పదవ} \te{సంవత్సరం}} (\emph{padava saṁvatsaraṁ}) — the tenth year.

That is the ten. Say them in order — \emph{modaṭi, reṇḍava, mūḍava, nālugava, aidava, ārava, ēḍava, enimidava, tommidava, padava} — and count what you actually had to learn: \textbf{one word and one ending}. Latin needed ten separate words for the same stretch, and English still does.

\subsection*{Your turn}
Say these aloud:

\begin{itemize}
\raggedright
  \item \emph{padava}
  \item the ten in order, \emph{modaṭi} to \emph{padava}
  \item \emph{padava saṁvatsaraṁ} — the tenth year
\end{itemize}

\begin{tcolorbox}[breakable,colback=teal!4,colframe=teal!35!black,title={Before you move on}]
Say the ten in order. (\emph{\textbf{Modaṭi}} through \emph{\textbf{padava}}.) How much had to be learned as separate words? (\textbf{One} — \emph{modaṭi}.) And say \textquotedblleft{}the tenth year\textquotedblright{}. (\textbf{\te{పదవ} \te{సంవత్సరం}}.)
\end{tcolorbox}
\section[padakoṇḍava]{\te{పదకొండవ} (padakoṇḍava) — the ending does not stop at ten}

\label{lesson:TE-C82-eleventh}



\begin{quote}
Recalls first, and then the last thing the chapter has to give.

\begin{itemize}
\raggedright
  \item \emph{Recall:} say \emph{padava} — \textbf{R1}, one lesson back, before it decays
  \item \emph{Recall:} say \emph{ārava} — \textbf{R2}, the first real retrieval, five to fifteen lessons back
  \item \emph{Recall:} say \emph{māma} — \textbf{R3}, durable, twenty to sixty lessons back
  \item \emph{Recall:} say \emph{pulupu} — \textbf{R3}, a second word from the same distance
  \item \emph{Recall:} say \emph{ekkaḍa} — \textbf{R4}, recognition at distance, eighty lessons back or more
  \item \emph{Recall:} say \emph{the i- / a- / e- pattern} — \textbf{R4}, a second word from the same distance
\end{itemize}
\end{quote}

\begin{grammarlens}[title={Grammar Lens: every number you have, and every one you meet}]
\textbf{\te{పదకొండవ}} (\emph{padakoṇḍava}) — \textquotedblleft{}\textbf{eleventh}\textquotedblright{}.

\textbf{\te{పదకొండు}} $\to$ \textbf{\te{పదకొండవ}}. Nothing new: the same final vowel goes, the same \textbf{-\te{వ}} arrives.

That is the point of this last lesson. The ending was never a fact about the numbers one to ten; it is a fact about \textbf{numbers}. Take any Telugu number you know and it works:

\noindent
\begin{tabularx}{\linewidth}{@{}>{\raggedright\arraybackslash}X>{\raggedright\arraybackslash}X@{}}
\toprule
\textbf{number} & \textbf{ordinal} \\
\midrule
\textbf{\te{పదకొండు}} (11) & \textbf{\te{పదకొండవ}} \\
\textbf{\te{ఇరవై}} (20) & \textbf{\te{ఇరవయ్యవ}} \\
\bottomrule
\end{tabularx}

So there is no second list to learn above ten, and no list to learn below it either except the single word \textbf{\te{మొదటి}}. One irregular word, one ending, and every position on the number line is sayable.
\end{grammarlens}

\subsection*{Your turn}
Say these aloud:

\begin{itemize}
\raggedright
  \item \emph{padakoṇḍava}
  \item the swap on a bigger number — \emph{padakoṇḍu} $\to$ \emph{padakoṇḍava}
  \item \emph{modaṭi} is the only word; \emph{-va} is the rest
\end{itemize}

\begin{tcolorbox}[breakable,colback=teal!4,colframe=teal!35!black,title={Before you move on}]
Say the eleventh. (\textbf{\te{పదకొండవ}}.) How many separate ordinal words does Telugu have to be learned? (\textbf{One} — \emph{modaṭi}.) And what makes all the others? (\textbf{The number, minus its final vowel, plus -\te{వ}}.)
\end{tcolorbox}
